% Part: first-order-logic
% Chapter: axiomatic-deduction
% Section: provability-quantifiers

% verification of properties of provability needed for maximally
% consistent sets in the completeness chapter.

\documentclass[../../../include/open-logic-section]{subfiles}

\begin{document}

\olfileid[fr]{fol}{axd}{qpr}

\olsection{\usetoken{S}{derivability} et quantificateurs}

\begin{explain}
La preuve du théorème de complétude exige également que les
!!{derivation}s axiomatiques établissent les propriétés de~$\Proves$
présentées dans cette section.
\end{explain}

\begin{thm}
\ollabel{thm:strong-generalization}
Si $c$ est !!a{constant} qui ne figure ni dans~$\Gamma$ ni dans
$!A(x)$, et si $\Gamma \Proves !A(c)$, alors
$\Gamma \Proves \lforall[x][!A(x)]$.
\end{thm}

\begin{proof}
Par le théorème de la déduction,
$\Gamma \Proves \ltrue \lif !A(c)$. Comme $c$ ne figure ni dans
$\Gamma$ ni dans~$\ltrue$, la règle~\QR{} donne
$\Gamma \Proves \ltrue \lif \lforall[x][!A(x)]$. Enfin,
$\Gamma \Proves \ltrue$ puisque $\ltrue$ est un axiome; le modus
ponens entraîne donc $\Gamma \Proves \lforall[x][!A(x)]$.
\end{proof}

\begin{prop}
\ollabel{prop:provability-quantifiers}
\begin{tagenumerate}{prvEx,prvAll}
\tagitem{prvEx}{$!A(t) \Proves \lexists[x][!A(x)]$.}{}

\tagitem{prvAll}{$\lforall[x][!A(x)] \Proves !A(t)$.}{}
\end{tagenumerate}
\end{prop}

\begin{proof}
\begin{tagenumerate}{prvEx,prvAll}
\tagitem{prvEx}{D'après l'\olref[qua]{ax:q2} et le théorème de la
déduction.}{}
\tagitem{prvAll}{D'après l'\olref[qua]{ax:q1} et le théorème de la
déduction.}{}
\end{tagenumerate}
\end{proof}

\begin{editorial}
À la fin de la première preuve, la source invoque une seconde fois le
théorème de la déduction pour passer de
$\Gamma \Proves \ltrue \lif \lforall[x][!A(x)]$ à
$\Gamma \Proves \lforall[x][!A(x)]$. Ce passage requiert en réalité
l'axiome~$\ltrue$ et le modus ponens; ces deux étapes sont explicitées.
\end{editorial}

\end{document}
